\newcommand{\figuraGravitacionContinua}{%
\begin{figure}[H]
    \centering
    \begin{tikzpicture}[>=Latex,font=\small]
        \node[font=\bfseries] at (6.1,5.55)
            {Campo gravitacional de una distribución continua};

        \draw[very thick,fill=blue!8]
            plot[smooth cycle,tension=0.75] coordinates {
                (0.65,1.05) (1.25,3.80) (3.20,4.70)
                (5.15,3.85) (5.55,1.55) (3.75,0.55)};
        \node[blue!65!black] at (2.95,4.10) {$\rho(\vec{r}\,')$};

        \foreach \x/\y in {
            1.25/1.45,1.55/2.55,1.90/3.45,2.35/1.10,
            2.55/2.30,2.90/3.20,3.30/1.55,3.65/2.65,
            4.10/3.55,4.55/1.70,4.75/2.70}
        {
            \shade[ball color=blue!60] (\x,\y) circle (0.075);
        }

        \coordinate (S) at (3.65,2.65);
        \coordinate (P) at (9.45,3.55);
        \fill[red!75!black] (S) circle (0.10);
        \fill[black] (P) circle (0.08);
        \node[anchor=north,fill=blue!4,inner sep=2pt]
            (dm) at (3.10,1.25)
            {$dm=\rho(\vec{r}\,')\,d^3\vec{r}\,'$};
        \draw[gray!70,thin] (dm.north east)--(S);
        \node[anchor=south west] at ($(P)+(0.08,0.08)$)
            {punto de campo $\vec{r}$};

        \draw[->,very thick,red!75!black]
            (S)--(P)
            node[midway,above,sloped]
            {$\vec{r}-\vec{r}\,'$};
        \draw[->,ultra thick,green!45!black]
            (P)--++(-1.55,-0.25)
            node[midway,below] {$d\vec{g}$};

        \draw[->,thick] (0.20,0.25)--(1.25,0.25)
            node[below] {$x$};
        \draw[->,thick] (0.20,0.25)--(0.20,1.30)
            node[left] {$y$};

        \node[align=center,text=black!70] at (8.85,2.15)
            {superposición de las contribuciones\\
             de todos los elementos de masa};
    \end{tikzpicture}
    \caption{Cada elemento de masa de la distribución produce una
    contribución diferencial al campo gravitacional. El campo total se
    obtiene por superposición, integrando sobre todas las posiciones fuente
    $\vec{r}\,'$.}
    \label{fig:gravitacion-continua}
\end{figure}%
}